\documentclass[ 12pt, a4paper]{article}
\usepackage[ ]{indentfirst}

\title{Revisão P1\\Construção de Compiladores\\Prof. Bráulio de Mello}
\author{Erickson G. Müller}
\date{23 de Abril de 2026}

\begin{document}
\maketitle
\section*{Conteúdos}
	\begin{itemize}
		\item Construção de Analisadores Léxicos
		\item Reconhecimento Sintático por APND
		\item APND, Analisador Sintático Preditivo
		\item Analisadores Left-Right
			\begin{itemize}
				\item Construção de tabela de parsing SLR
				\item Algoritmo de mapeamento do parser LR para análise sintática.
			\end{itemize}
	\end{itemize}
\newpage


\section{Construção de Analisador Léxico}

	A saída do Analisador Léxico é a entrada do Analisador Sintático.

\textbf{Processo:}
	\begin{enumerate}
		\item Definir tokens
		\item Construir o AFND
		\item Determinar o AF
		\item Acrescentar estado de erro para tratar tokens inválidos
		\item Construir o analisador (algoritmo)
		\item Tratamento de erro
		\begin{enumerate}
			\item Rótulo não reconhecido
			\item Descrição do erro
			\item Adicionar na Tabela de Símbolos
		\end{enumerate}
		\item Reconhecimento
	\end{enumerate}
	
	1 estado reconhece 1 token
	
	1 Token pode ser reconhecido por mais de 1 estado

	\subsection{Tabela de Símbolos}
		A tabela de símbolos segue a estrutura:
		
		\textit{Fita: [sequência de estados]+\$}
		
		\textit{L1: rótulo da cadeia}
		
		\textit{[Mensagem de erro]}
\end{document}